% SEGMENT OLP-0099-S01
% Part: first-order-logic
% Chapter: tableaux
% Section: rules-and-proofs

\documentclass[../../../include/open-logic-section]{subfiles}

\begin{document}

\iftag{FOL}
      {\olfileid{fol}{tab}{rul}}
      {\olfileid{pl}{tab}{rul}}

\olsection{விதிகளும் \usetoken{P}{tableau}களும்}

!!^a{tableau} என்பது !!a{sentence} !!a{structure}வில் மெய்யாகவோ
பொய்யாகவோ அமையக்கூடிய வழிகளை முறையாக ஆய்வதாகும். ஓர் அட்டவணை மரத்தின்
கட்டுமானக் கூறுகள் !!{signed formula}s: ஒரு !!{sentence} உடன்
$\True$ அல்லது~$\False$ என்ற மெய்மதிப்புக் ``குறி'' ஒன்றை இணைத்தவை. இந்தக்
குறியிட்ட !!{formula}s கீழ்நோக்கி வளரும் மரவுரு ஒன்றில் ஒழுங்கமைக்கப்படுகின்றன.

% SEGMENT OLP-0099-S02
\begin{defn}
  \emph{!!^a{signed formula}} என்பது ஒரு மெய்மதிப்பும் !!a{sentence} உம்
  கொண்ட சோடி; அதாவது பின்வரும் இரண்டில் ஒன்று:
  \[
  \sFmla{\True}{!A} \text{ அல்லது } \sFmla{\False}{!A}.
  \]
\end{defn}

% SEGMENT OLP-0099-S03
உள்ளுணர்வு நோக்கில், $\sFmla{\True}{!A}$ என்பதை ``$!A$ மெய்யாக இருக்கலாம்''
என்றும், $\sFmla{\False}{!A}$ என்பதை ``$!A$ பொய்யாக இருக்கலாம்'' என்றும்
(ஏதோவொரு !!{structure}வில்) வாசிக்கலாம்.

% SEGMENT OLP-0099-S04
மரவுருவிலுள்ள ஒவ்வொரு !!{signed formula} உம் ஓர் \emph{எடுகோள்}
(மரவுருவின் மிக மேலே பட்டியலிடப்படும்) அல்லது அதற்கு மேலுள்ள
!!a{signed formula}விலிருந்து பல அனுமான விதிகளில் ஒன்றால் பெறப்படுவது. முந்தைய
!!{formula}வின் ஒவ்வொரு சாத்தியமான !!{main operator}க்கும் இரண்டு விதிகள்
உள்ளன: குறி~$\True$ ஆகும் நேர்வுக்கு ஒன்றும், குறி~$\False$ ஆகும் நேர்வுக்கு
ஒன்றும். சில விதிகள் மரவுருவைக் கிளைக்கச் செய்கின்றன; மற்றவை கிளையில்
!!{signed formula}s ஐ மட்டும் சேர்க்கின்றன. ஒரு விதியை அதற்கு உடனடியாக
முந்திய !!{signed formula}விற்கே பயன்படுத்த வேண்டியதில்லை; வேரிலிருந்து விதி
பயன்படுத்தப்படும் இடம்வரை உள்ள கிளையின் எந்த !!{signed formula}விற்கும் அது
பயன்படுத்தப்படலாம் (பெரும்பாலும் பயன்படுத்தப்பட வேண்டும்).

% SEGMENT OLP-0099-S05
ஒரு கிளையில் $\sFmla{\True}{!A}$, $\sFmla{\False}{!A}$ ஆகிய இரண்டும்
இருந்தால் அது \emph{மூடியது}. ஒவ்வொரு கிளையும் மூடியுள்ள !!a{tableau}
மூடிய அட்டவணை மரமாகும். உள்ளுணர்வு விளக்கத்தின்படி, ஒவ்வொரு கிளையும்
ஒருங்கே நிகழக்கூடிய வாய்ப்பு ஒன்றை விவரிக்கிறது; ஆனால்
$\sFmla{\True}{!A}$ உம் $\sFmla{\False}{!A}$ உம் ஒருங்கே நிகழ முடியாது.
வேறு சொற்களில், ஒரு கிளை மூடியிருந்தால் அது விவரிக்கும் வாய்ப்பு நீக்கப்பட்டு
விட்டது. குறிப்பாக, $\sFmla{\True}{!A}$ வடிவிலுள்ள ஒவ்வோர் எடுகோளையும்
மெய்யாக்கியும் $\sFmla{\False}{!A}$ வடிவிலுள்ள ஒவ்வோர் எடுகோளையும்
பொய்யாக்கியும் ஒரே நேரத்தில் அமையும் அனைத்து வாய்ப்புகளையும் ஒரு மூடிய
!!{tableau} நீக்குகிறது.

% SEGMENT OLP-0099-S06
$!A$ \emph{க்கான மூடிய !!a{tableau}} என்பது
$\sFmla{\False}{!A}$ ஐ வேராகக் கொண்ட மூடிய !!a{tableau}. அத்தகைய மூடிய
!!a{tableau} இருந்தால், $!A$ பொய்யாக இருக்கக்கூடிய எல்லா வாய்ப்புகளும்
நீக்கப்பட்டுள்ளன; அதாவது ஒவ்வொரு !!{structure}விலும் $!A$ மெய்யாகவே இருக்க
வேண்டும்.

\end{document}
